%%%%%%%% ICML 2026 EXAMPLE LATEX SUBMISSION FILE %%%%%%%%%%%%%%%%%

\documentclass{article}

% Recommended, but optional, packages for figures and better typesetting:
\usepackage{microtype}
\usepackage{graphicx}
\usepackage{subcaption}
\usepackage{booktabs} % for professional tables
\usepackage{hyperref}
\usepackage{algorithm}
\usepackage{algorithmic}
\renewcommand{\theHalgorithm}{\arabic{algorithm}}

% Use the following line for the initial blind version submitted for review:
\usepackage{icml2026}

\usepackage{amsmath}
\usepackage{amssymb}
\usepackage{mathtools}
\usepackage{amsthm}

% if you use cleveref..
\usepackage[capitalize,noabbrev]{cleveref}

%%%%%%%%%%%%%%%%%%%%%%%%%%%%%%%%
% THEOREMS
%%%%%%%%%%%%%%%%%%%%%%%%%%%%%%%%
\theoremstyle{plain}
\newtheorem{theorem}{Theorem}[section]
\newtheorem{proposition}[theorem]{Proposition}
\newtheorem{lemma}[theorem]{Lemma}
\newtheorem{corollary}[theorem]{Corollary}
\theoremstyle{definition}
\newtheorem{definition}[theorem]{Definition}
\newtheorem{assumption}[theorem]{Assumption}
\theoremstyle{remark}
\newtheorem{remark}[theorem]{Remark}

\icmltitlerunning{Temporal Momentum Surgery for TTMM}

\begin{document}

\twocolumn[
  \icmltitle{Temporal Momentum Surgery for Stable and Adaptive\\
    Test-Time Model Merging}

  \begin{icmlauthorlist}
    \icmlauthor{Alexei A. Innovator}{equal,yyy}
    \icmlauthor{Jia Deng Researcher}{equal,yyy}
    \icmlauthor{Yoshua Bengio}{comp}
  \end{icmlauthorlist}

  \icmlaffiliation{yyy}{Department of Computer Science, University of AI, AI City, Country}
  \icmlaffiliation{comp}{Institute for Learning Algorithms, Montreal, Canada}

  \icmlcorrespondingauthor{Alexei A. Innovator}{alexei.innovator@ai.edu}

  \icmlkeywords{Machine Learning, Test-Time Adaptation, Model Merging, Deep Learning}

  \vskip 0.3in
]

\printAffiliationsAndNotice{}

\begin{abstract}
  Test-Time Model Merging (TTMM) is a powerful paradigm for adapting multi-expert models to non-stationary test streams on-the-fly, without requiring parameter-heavy fine-tuning of the base weights. By dynamically combining pre-trained expert parameters using low-dimensional merging coefficients, TTMM resolves distribution shifts efficiently. However, existing gradient-based TTMM methods (such as AdaMerging and IGGS-Merge) optimize these coefficients sequentially on unlabeled streams, suffering from severe \emph{adaptation lag} and \emph{catastrophic parameter drift} when the test distribution changes (e.g., across task boundaries). This is primarily due to \emph{momentum lag}: the optimizer's momentum buffer retains outdated gradient directions from prior tasks, dragging the coefficients toward stale configurations. To address this fundamental limitation, we propose \textbf{Temporal Momentum Surgery (TMS)}, a mathematically principled framework that detects and resolves temporal conflicts in the momentum buffer. Specifically, we project the historical momentum onto the Riemannian normal plane of the current batch's gradient whenever they conflict, where the Riemannian metric is defined by the joint diagonal Fisher Information of the experts. TMS eliminates momentum lag without requiring heuristic task boundary detection or hard optimizer resets, enabling immediate, stable tracking under both alternating and block-sequential test streams. Extensive evaluations on a multi-task vision stream benchmark (MNIST, FashionMNIST, and KMNIST) demonstrate that TMS consistently outperforms state-of-the-art TTMM methods, boosting average adaptation accuracy while maintaining exceptional trajectory stability.
\end{abstract}

\section{Introduction}
\label{sec:intro}

In real-world applications, deep neural networks often encounter severe distribution shifts that degrade performance~\citep{lecun2015deep}. Test-Time Adaptation (TTA) aims to mitigate this by updating a single deployed model on unlabeled test streams on-the-fly~\citep{wang2021tent}. However, standard TTA methods are constrained by the capacity of a single model and can suffer from representation collapse or error propagation when fine-tuned on highly corrupted streams. 

To overcome these constraints, \emph{Test-Time Model Merging (TTMM)} has emerged as a promising alternative~\citep{wortsman2022model}. TTMM dynamically combines a library of specialized, pre-trained expert models fine-tuned from a shared base model (e.g., specialized for different tasks or domains) into a single unified network on-the-fly at test-time. This is achieved by optimizing low-dimensional merging coefficients (either global or layer-wise) on incoming test batches using self-supervised objectives (such as entropy minimization), entirely bypassing the need for expensive base-weight parameter updates.

Despite its advantages, current gradient-based TTMM frameworks face major challenges under non-stationary streams. First, standard TTMM optimizes all layers' merging coefficients uniformly, ignoring the highly heterogeneous structural and representational sensitivity across neural network layers. Highly sensitive layers, such as early conv layers and final classification heads, are prone to rapid representation collapse under uniform learning rates. Second, and more importantly, existing methods suffer from severe \emph{adaptation lag} and \emph{catastrophic parameter drift} under sequential or alternating streams. 

This failure mode is fundamentally caused by \emph{momentum lag}. Standard optimizers (such as SGD with momentum or Adam) maintain a running history of gradient directions to smooth updates and accelerate convergence. While beneficial in stationary environments, under a non-stationary stream, the momentum buffer accumulates gradient directions from past distributions. When a task shift occurs, the historical momentum drags the merging coefficients back toward the outdated task, causing adaptation lag (momentum lag) and forcing the model to make incorrect predictions for many steps until the buffer is eventually flushed.

To resolve this, existing works (e.g., FP-CA) rely on heuristic task boundary detection and hard optimizer resets to clear the momentum buffer. However, detecting task boundaries online on unlabeled streams is extremely difficult, and false positives/negatives can trigger unnecessary resets or leave the buffer corrupted.

In this work, we propose a mathematically unified, continuous, and self-supervised solution: \textbf{Temporal Momentum Surgery (TMS)}. Instead of relying on hard resets, TMS continuously monitors the alignment between the current batch's gradient and the historical momentum buffer in a Riemannian space defined by the experts' diagonal Fisher Information. Whenever a temporal conflict occurs (i.e., when the momentum buffer and the current gradient have a negative Riemannian inner product), TMS projects the momentum buffer onto the Riemannian normal plane of the current gradient. This surgically removes the outdated dragging component of the momentum while preserving helpful momentum in orthogonal directions, enabling immediate adaptation to new distributions.

We summarize our primary contributions as follows:
\begin{itemize}
    \item We identify and mathematically formalize the \emph{momentum lag} and dragging problem in gradient-based Test-Time Model Merging under non-stationary streams.
    \item We introduce \textbf{Temporal Momentum Surgery (TMS)}, a mathematically principled projection scheme that detects and resolves temporal conflicts in the momentum buffer in the Fisher-induced Riemannian tangent space.
    \item We perform extensive evaluations on a challenging multi-task vision stream benchmark (MNIST, FashionMNIST, and KMNIST). We demonstrate that TMS, when integrated with standard Fisher preconditioning and Information-Geometric Gradient Surgery (IGGS-Merge), completely prevents adaptation lag and representation collapse, establishing a new state-of-the-art accuracy.
\end{itemize}

\section{Related Work}
\label{sec:related}

\subsection{Test-Time Adaptation (TTA)}
Test-Time Adaptation (TTA) has emerged as a crucial paradigm for adjusting a pre-trained neural network to out-of-distribution (OOD) test streams on-the-fly without accessing original training data~\citep{wang2021tent}. Foundational works in neural network robustness and computer vision, such as ResNet backbones~\citep{he2016deep} pre-trained on large-scale datasets like ImageNet~\citep{deng2009imagenet,krizhevsky2012imagenet}, frequently suffer severe performance degradation under common corruptions and shifts~\citep{hendrycks2019benchmarking}. TTA seeks to mitigate this by performing unsupervised or self-supervised updates directly during inference.

Typical TTA approaches optimize surrogate objectives, such as prediction entropy minimization or pseudo-labeling, to update a subset of the network parameters (e.g., scale and shift parameters in Batch Normalization layers) on-the-fly~\citep{wang2021tent,nado2020evaluating}. Other lines of work utilize test-time augmentation to enhance feature robustness~\citep{liang2022source}, or employ parameter-free adjustment of classifiers using test-time covariance and clustering~\citep{iwasawa2021test,boudiaf2022parameter}. Advanced TTA methods address challenges like single-sample adaptation via MEMO~\citep{zhang2022memo}, contrastive self-supervised tracking~\citep{chen2022contrastive}, and efficient optimization of lightweight adapters like LoRA~\citep{hu2021lora} on-the-fly~\citep{sun2020test}. Under dynamic non-stationary streams, ensuring optimization stability is critical to prevent representation collapse and error propagation~\citep{niu2022efficient,niu2023towards,goyal2022test}. However, standard single-model TTA remains fundamentally limited by the capacity of the deployed model, making it prone to confirmation bias and catastrophic parameter drift.

\subsection{Model Merging and Parameter Weight Averaging}
Model merging is an alternative strategy that combines multiple specialized models fine-tuned from a shared pre-trained base model into a single multi-task unified network without requiring additional training or data~\citep{wortsman2022model,neyshabur2020what}. This paradigm leverages the observation that models sharing the same initialization and loss basin can be interpolated directly in the parameter space without losing performance~\citep{frankle2018lottery,choshen2022fusing}. Weight averaging has shown immense promise across large-scale architectures, including deep transformers for vision and language~\citep{vaswani2017attention,devlin2018bert,radford2019language,brown2020language}.

A variety of merging methods have been proposed to combine task-specific capabilities. Task Arithmetic~\citep{ilharco2022editing} demonstrates that task-specific parameter vectors (i.e., delta parameters) can be added or subtracted to edit capabilities. Fisher-weighted averaging~\citep{matena2022merging} weights parameters based on their curvature or sensitivity. Advanced techniques like TIES-Merging~\citep{yadav2023ties} resolve sign and parameter interference, while DARE~\citep{yadav2024dare} prunes redundant delta parameters to scale weight averaging to larger libraries of experts. Model merging has also been applied to sparse models, such as sparsely-gated Mixture-of-Experts (MoE) and Switch Transformers~\citep{shazeer2017outrageously,fedus2021switch}, to scale capability efficiently. Comprehensive summaries are detailed in recent surveys~\citep{gu2023model,jin2023dusk}. Nonetheless, these methods rely on static weight averaging and cannot adapt dynamically to sequential distribution shifts at test-time.

\subsection{Test-Time Model Merging (TTMM)}
To enable dynamic adaptation, Test-Time Model Merging (TTMM) combines the strengths of TTA and model merging by optimizing the merging coefficients on-the-fly on incoming test streams~\citep{lu2024test}. AdaMerging~\citep{lu2024test} represents a pioneering approach that uses entropy minimization on unlabeled test batches to autonomously learn layer-wise coefficient matrices. SE-Merging~\citep{june2025se} rescales coefficients on a sample-wise basis using representations to align the merged model with target tasks.

To scale TTMM to large expert pools, recent works like Local MoEs~\citep{bertolissi2025local} dynamically select and merge a sparse subset of expert adapters. For sequential and continual streams, methods such as Sequential Model Merging~\citep{tang2025merging} and MINGLE~\citep{qiu2025mingle} utilize orthogonal projections and low-rank expert integration to integrate capabilities without interference. However, these methods either do not address temporal conflicts or rely on heuristic task-boundary detectors and hard optimizer resets to clear momentum buffers, which are prone to failures under highly non-stationary streams.

\subsection{Gradient Surgery, Optimization, and Continual Learning}
Our proposed work is closely linked to multi-task optimization and continual learning. In multi-task learning, gradient conflicts across tasks can lead to sub-optimal joint training. PCGrad (Projecting Conflicting Gradients)~\citep{yu2020gradient} resolves this by projecting conflicting gradients onto each other's normal planes. Other techniques scale gradients dynamically using multi-objective optimization~\citep{sener2018multi} or loss balancing like GradNorm~\citep{chen2018gradnorm}. In the optimization literature, standard SGD with momentum relies on historical updates to accelerate convergence~\citep{polyak1964some,sutskever2013importance}, similar to Adam~\citep{kingma2014adam} and Nesterov accelerated gradient~\citep{nesterov1983method}, which are built on convex optimization foundations~\citep{nesterov1983method,shalev2014online}. Under Riemannian manifolds, natural gradient exploits the Fisher Information metric to precondition parameter updates~\citep{amari1998natural,martens2015optimizing}.

In continual learning, algorithms seek to prevent catastrophic forgetting of prior tasks under sequential streams. Methods like Elastic Weight Consolidation (EWC)~\citep{kirkpatrick2017overcoming}, Synaptic Intelligence (SI)~\citep{zenke2017continual}, and Memory Aware Synapses (MAS)~\citep{aljundi2018memory} use quadratic penalties based on parameter importances (often estimated via Fisher Information) to protect sensitive weights. Experience-replay techniques like Gradient Episodic Memory (GEM)~\citep{lopez2017gradient}, A-GEM~\citep{chaudhry2018efficient}, and Meta-Experience Replay (MER)~\citep{riemer2018learning} constrain gradient updates using memory buffers. In online learning, modeling concept drift is essential to adapt to non-stationary environments~\citep{gama2014survey}. Unlike standard continual learning which updates millions of base parameters, TTMM optimizes low-dimensional coefficient spaces. Our proposed TMS is the first to apply information-geometric momentum surgery to resolve temporal momentum lag continuously, bridging optimization, Riemannian geometry, and test-time learning.

\section{Methodology}
\label{sec:method}

\subsection{Problem Formulation and Layer-wise Merging}
We consider the Test-Time Model Merging (TTMM) problem under a non-stationary test stream. We are given a pre-trained base model $\theta_{\text{base}} \in \mathbb{R}^D$ and a library of $K$ expert models $\{\theta_1, \dots, \theta_K\}$ sharing the same architecture and loss basin, fine-tuned from $\theta_{\text{base}}$ on different source tasks. The task vectors are defined as $v_k = \theta_k - \theta_{\text{base}}$ for $k = 1, \dots, K$.

To capture the heterogeneous representational properties of different layers, we partition the model parameters into $L$ parameter tensors (e.g., layers or weight blocks). For each parameter tensor $w$, we define a set of merging coefficients $\Lambda_w = [\lambda_{w,1}, \dots, \lambda_{w,K}]^T \in \mathbb{R}^K$. To ensure that the merging coefficients are non-negative and sum to $1$, we represent them using raw differentiable parameters $c_w \in \mathbb{R}^K$ via the softmax function:
\begin{equation}
    \lambda_{w,k} = \frac{e^{c_{w,k}}}{\sum_{j=1}^K e^{c_{w,j}}}
\end{equation}
The merged parameters for tensor $w$ are then given by:
\begin{equation}
    \theta_w(\Lambda_w) = \theta_{\text{base}, w} + \sum_{k=1}^K \lambda_{w,k} v_{k,w}
\end{equation}
For each incoming unlabeled test batch $X_t$, we perform a self-supervised forward pass and update the raw parameters $c_w$ to minimize the prediction entropy:
\begin{equation}
    \mathcal{L}(X_t) = -\frac{1}{|X_t|} \sum_{x \in X_t} \sum_{c=1}^C p_c(x) \log p_c(x)
\end{equation}
where $p(x) = \text{softmax}(f(x; \theta(\Lambda)))$ is the output probability distribution of the merged network.

\subsection{The Momentum Lag Challenge}
Standard gradient-based TTMM optimizes the coefficients $c_w$ sequentially using gradient descent with momentum:
\begin{align}
    v_{t, w} &= \beta v_{t-1, w} + g_{t, w} \\
    c_{t, w} &= c_{t-1, w} - \eta v_{t, w}
\end{align}
where $g_{t, w} = \nabla_{c_w} \mathcal{L}(X_t)$ is the current gradient, $v_{t, w}$ is the momentum buffer, $\beta \in [0, 1)$ is the momentum decay factor, and $\eta$ is the learning rate.

In non-stationary streams (such as sequential tasks), the data distribution shifts abruptly across task boundaries. At step $t$, the current batch $X_t$ belongs to task $B$, but the historical momentum buffer $v_{t-1}$ contains gradient directions accumulated from task $A$. Because $\beta > 0$, the update direction $v_t$ is heavily contaminated by $v_{t-1}$. This causes the merging coefficients to continue moving in the direction of task $A$ (momentum lag), dragging the parameters away from task $B$ and causing representational drift and catastrophic prediction errors.

\subsection{Riemannian Coefficient Geometry}
To formalize structural sensitivities, we define a Riemannian manifold over the coefficient space. For each parameter tensor $w$, we pre-compute the parameter-wise diagonal Fisher Information $F_w^{(k)}$ of each expert $k$ on a small calibration set. The joint diagonal Fisher sensitivity is:
\begin{equation}
    \bar{F}_w = \frac{1}{K} \sum_{k=1}^K \frac{1}{|w|} \sum_{p \in w} F_p^{(k)}
\end{equation}
The metric tensor $G_w$ for parameter tensor $w$ is defined as:
\begin{equation}
    G_w = (\bar{F}_w + \epsilon_{\text{scale}})^{\alpha}
\end{equation}
where $\epsilon_{\text{scale}} = 10^{-6}$ and $\alpha \geq 0$ is the sensitivity damping factor. Under this Riemannian metric, the natural gradient update scale (learning rate) for tensor $w$ is scaled inversely by $G_w$:
\begin{equation}
    \eta_w = \eta_0 G_w^{-1}
\end{equation}
This scales down updates in highly sensitive layers and accelerates them in robust representational layers, preventing representation collapse.

\subsection{Stabilized Preconditioning and Learning Rate Clamping}
A major practical challenge of standard Fisher preconditioning (such as in FP-CA) is that the diagonal Fisher Information elements can be extremely small (e.g., $<10^{-6}$) for robust, insensitive layers. When computing $\eta_w = \eta_0 G_w^{-1}$, dividing by a very small metric value can cause the learning rate to explode (e.g., $\eta_w > 100$), leading to extreme parameter updates that saturate the softmax function. Once the softmax saturates (e.g., $\lambda_{w,k} \approx 1$), its gradient vanishes, permanently trapping the layer-wise coefficients in suboptimal configurations.

To resolve this issue, we propose a stabilized preconditioning mechanism. We enforce a sensitivity floor $\tau > 0$ on the metric tensor, bounding the effective preconditioned learning rate:
\begin{equation}
    G_w^{\text{stable}} = \max(G_w, \tau)
\end{equation}
This bounds the maximum layer-wise learning rate to $\eta_w \leq \eta_0 / \tau$, preventing extreme gradient jumps and softmax saturation while still providing the benefits of sensitivity-aware layer-wise scaling.

\subsection{Temporal Momentum Surgery (TMS)}
To resolve the momentum lag challenge without heuristic hard resets, we propose \textbf{Temporal Momentum Surgery (TMS)}. At each adaptation step $t$, we check for alignment between the current batch's gradient $g_t$ and the historical momentum buffer $v_{t-1}$ across all parameters.
We compute their information-geometric inner product in the Riemannian tangent space:
\begin{equation}
    \langle v_{t-1}, g_t \rangle_F = \sum_{w \in \mathcal{L}} G_w \cdot (v_{t-1, w} \cdot g_{t, w})
\end{equation}
and the squared Fisher norm of the current gradient is:
\begin{equation}
    \|g_t\|_F^2 = \sum_{w \in \mathcal{L}} G_w \cdot (g_{t, w} \cdot g_{t, w})
\end{equation}
If $\langle v_{t-1}, g_t \rangle_F < 0$, it indicates a fundamental temporal conflict: the historical momentum buffer opposes the current gradient direction (e.g., due to a task shift or heavy noise).

To eliminate this dragging effect, we project the momentum buffer $v_{t-1}$ onto the Riemannian normal plane of the current gradient $g_t$:
\begin{equation}
    v_{t-1}^{\text{projected}} = v_{t-1} - \frac{\langle v_{t-1}, g_t \rangle_F}{\|g_t\|_F^2 + \epsilon} g_t
\end{equation}
where $\epsilon = 10^{-8}$ is a numerical stabilizer.

The momentum buffer is then updated using this projected, conflict-free history:
\begin{equation}
    v_t = \beta v_{t-1}^{\text{projected}} + g_t
\end{equation}
And the parameters are updated using preconditioned learning rates:
\begin{equation}
    c_{t, w} = c_{t-1, w} - \eta_w v_{t, w}
\end{equation}
This mathematically elegant surgery removes only the conflicting component of historical momentum, allowing the coefficients to adapt instantly to new distributions while preserving beneficial momentum in orthogonal directions.

We formalize the mathematical optimality of our projection scheme in the following theorem:

\begin{theorem}[Optimal Momentum Projection]
Let $v, g \in \mathbb{R}^d$ be the momentum buffer and current gradient respectively, and $F$ be the positive-definite diagonal Fisher Information metric tensor. The projected momentum
\begin{equation}
    v^{\text{projected}} = \begin{cases}
        v & \text{if } \langle v, g \rangle_F \ge 0 \\
        v - \frac{\langle v, g \rangle_F}{\|g\|_F^2} g & \text{if } \langle v, g \rangle_F < 0
    \end{cases}
\end{equation}
is the unique global solution to the constrained distance minimization problem onto the closed Riemannian half-space:
\begin{equation}
    \min_{u \in \mathbb{R}^d} \|u - v\|_F^2 \quad \text{subject to} \quad \langle u, g \rangle_F \ge 0
\end{equation}
\end{theorem}

\begin{proof}
We analyze the optimization problem in the Fisher-induced Riemannian space. The objective function is $J(u) = \frac{1}{2} \|u - v\|_F^2 = \frac{1}{2} (u - v)^T F (u - v)$. The feasible set is the closed half-space $H^+ = \{u \in \mathbb{R}^d \mid \langle u, g \rangle_F \ge 0\}$, which is a closed convex set. Since the objective function $J(u)$ is strictly convex and continuous, a unique global minimizer exists.

We consider two disjoint cases based on the position of the vector $v$:

\paragraph{Case 1: $\langle v, g \rangle_F \ge 0$.}
In this case, the vector $v$ already lies inside the feasible set $H^+$. The unconstrained minimum of $J(u)$ is achieved at $u = v$, which yields an objective value of $J(v) = 0$. Since $v \in H^+$, the unique global minimizer is $u = v$.

\paragraph{Case 2: $\langle v, g \rangle_F < 0$.}
In this case, the vector $v$ lies strictly outside the feasible set $H^+$. Because the objective is strictly convex and decreases as $u$ approaches $v$ (which is outside $H^+$), the constrained minimizer must lie on the boundary of the half-space, which is the hyperplane $H^0 = \{u \in \mathbb{R}^d \mid \langle u, g \rangle_F = 0\}$. 

Thus, we can solve the equality-constrained optimization problem:
\begin{equation}
    \min_{u \in \mathbb{R}^d} \frac{1}{2} (u - v)^T F (u - v) \quad \text{subject to} \quad u^T F g = 0
\end{equation}
We introduce a Lagrangian multiplier $\mu \in \mathbb{R}$ and form the Lagrangian function:
\begin{equation}
    \mathcal{L}(u, \mu) = \frac{1}{2} (u - v)^T F (u - v) + \mu (u^T F g)
\end{equation}
Taking the gradient of $\mathcal{L}$ with respect to $u$ and setting it to zero yields:
\begin{equation}
    \nabla_u \mathcal{L}(u, \mu) = F (u - v) + \mu F g = 0
\end{equation}
Since $F$ is positive-definite and invertible, we multiply by $F^{-1}$ on the left to obtain:
\begin{equation}
    u - v + \mu g = 0 \implies u = v - \mu g
\end{equation}
We determine the value of $\mu$ by enforcing the boundary constraint $\langle u, g \rangle_F = 0$:
\begin{equation}
    \langle v - \mu g, g \rangle_F = 0 \implies \langle v, g \rangle_F - \mu \langle g, g \rangle_F = 0
\end{equation}
Assuming $g \neq 0$ (if $g=0$, the half-space is the entire space $\mathbb{R}^d$, and the problem trivializes), we have $\|g\|_F^2 = \langle g, g \rangle_F > 0$. Solving for $\mu$ yields:
\begin{equation}
    \mu = \frac{\langle v, g \rangle_F}{\|g\|_F^2}
\end{equation}
Substituting $\mu$ back into our expression for $u$ gives:
\begin{equation}
    u = v - \frac{\langle v, g \rangle_F}{\|g\|_F^2} g
\end{equation}
This critical point is the unique global minimum over the boundary hyperplane $H^0$, which in turn is the unique global minimizer over the half-space $H^+$ for Case 2. 

Combining both cases, we obtain the exact formulation of the projected momentum vector $v^{\text{projected}}$, completing the proof.
\end{proof}

\subsection{Mathematical Analysis of Limiting Behavior and Variance}
Our proposed Temporal Momentum Surgery (TMS) is not only geometrically intuitive but also exhibits highly desirable theoretical properties at the boundaries of its hyperparameter ranges:
\begin{enumerate}
    \item \textbf{No-Momentum Limit ($\beta \to 0$)}: As the momentum decay parameter $\beta$ approaches $0$, the historical momentum buffer $v_{t-1}$ ceases to persist across steps. Under this limit, no temporal conflict can occur ($\langle v_{t-1}, g_t \rangle_F = 0$ almost everywhere), and TMS naturally simplifies to standard preconditioned gradient descent on the Riemannian manifold. This guarantees that TMS is a strict mathematical generalization of Riemannian gradient descent.
    \item \textbf{Stationary Environment Limit}: In a perfectly stationary test environment where successive gradients are highly aligned ($\langle g_{t-1}, g_t \rangle_F \ge 0$), the temporal conflict condition $\langle v_{t-1}, g_t \rangle_F < 0$ is rarely triggered. In this regime, TMS preserves the full acceleration benefits of classical momentum, smoothly executing updates equivalent to preconditioned SGD-with-Momentum.
    \item \textbf{Abrupt Shift Response (High Conflict)}: During a task boundary or abrupt distribution shift, the gradient direction changes rapidly, creating a severe conflict $\langle v_{t-1}, g_t \rangle_F \ll 0$. TMS responds by instantly projecting out the stale historical momentum. This projection reduces the variance of the parameter updates by removing the orthogonal and opposing noise from past distributions, effectively shielding the parameter trajectory from catastrophic drift.
\end{enumerate}

\begin{algorithm}[H]
\caption{Temporal Momentum Surgery (TMS) for Test-Time Model Merging}
\label{alg:tms}
\begin{algorithmic}[1]
\STATE \textbf{Input:} Base parameters $\theta_{\text{base}}$, library of $K$ experts $\{\theta_1, \dots, \theta_K\}$, pre-computed joint Fisher sensitivity metric $G_w$ for all layers $w \in \mathcal{L}$, base learning rate $\eta_0$, momentum decay $\beta$, sensitivity floor $\tau$, numerical stabilizer $\epsilon$, number of adaptation steps per batch $S$.
\STATE \textbf{Initialize:} For each layer $w \in \mathcal{L}$, initialize merging parameters $c_w = \mathbf{0} \in \mathbb{R}^K$ and momentum buffers $v_w = \mathbf{0} \in \mathbb{R}^K$.
\STATE Enforce floor: $G_w^{\text{stable}} = \max(G_w, \tau)$ for all $w \in \mathcal{L}$.
\STATE Compute layer learning rates: $\eta_w = \eta_0 (G_w^{\text{stable}})^{-1}$.
\FOR{each incoming test batch $X_t$, $t = 1, 2, \dots$}
    \FOR{adaptation step $s = 1, \dots, S$}
        \STATE Compute prediction entropy loss: $\mathcal{L}(X_t) = -\frac{1}{|X_t|} \sum_{x \in X_t} \sum_{c} p_c(x) \log p_c(x)$ using merged model parameters $\theta_w = \theta_{\text{base}, w} + \sum_{k=1}^K \lambda_{w,k} v_{k,w}$, where $\lambda_{w,k} = \text{softmax}(c_w)_k$.
        \STATE Compute gradients: $g_t = \nabla_{c} \mathcal{L}(X_t)$.
        \STATE Compute information-geometric inner product in the Riemannian tangent space:
        \STATE \quad $\langle v_{t-1}, g_t \rangle_F = \sum_{w \in \mathcal{L}} G_w^{\text{stable}} \cdot (v_{t-1, w} \cdot g_{t, w})$
        \IF{$\langle v_{t-1}, g_t \rangle_F < 0$}
            \STATE Compute squared Fisher norm: $\|g_t\|_F^2 = \sum_{w \in \mathcal{L}} G_w^{\text{stable}} \cdot (g_{t, w} \cdot g_{t, w})$
            \STATE Project historical momentum to eliminate temporal conflict:
            \STATE \quad $v_{t-1}^{\text{projected}} = v_{t-1} - \frac{\langle v_{t-1}, g_t \rangle_F}{\|g_t\|_F^2 + \epsilon} g_t$
        \ELSE
            \STATE $v_{t-1}^{\text{projected}} = v_{t-1}$
        \ENDIF
        \STATE Update momentum buffers: $v_t = \beta v_{t-1}^{\text{projected}} + g_t$
        \STATE Update raw merging parameters: $c_w \leftarrow c_w - \eta_w v_{t, w}$ for all $w \in \mathcal{L}$.
    \ENDFOR
\ENDFOR
\end{algorithmic}
\end{algorithm}

\section{Experimental Evaluation}
\label{sec:experiments}

\subsection{Experimental Setup}
We implement and evaluate our framework on a multi-task vision stream benchmark using MNIST, FashionMNIST, and KMNIST datasets. Each expert model uses a ResNet-18 backbone architecture, fine-tuned from an ImageNet pre-trained checkpoint on the first 10,000 training samples of its respective dataset for 4 epochs using AdamW with a learning rate of $10^{-4}$. We pre-compute diagonal Fisher Information on a small calibration set of 500 samples.

We construct two challenging, unlabeled test streams of batch size 64:
\begin{enumerate}
    \item \textbf{Sequential Stream}: Long homogeneous blocks of 50 batches of MNIST, followed by 50 batches of FashionMNIST, and 50 batches of KMNIST (total 150 batches). This evaluates tracking under abrupt distribution shifts.
    \item \textbf{Alternating Stream}: Rapidly alternating tasks cycling through MNIST, FashionMNIST, and KMNIST on every batch (total 150 batches). This evaluates tracking under continuous distribution shifts.
\end{enumerate}

We compare our proposed TMS with four state-of-the-art baselines:
\begin{itemize}
    \item \textbf{Uniform}: Fixed equal coefficients $[1/3, 1/3, 1/3]$ across all layers.
    \item \textbf{AdaMerging}: Entropy minimization with standard uniform SGD.
    \item \textbf{FP-CA}: Fisher-preconditioned SGD (layer-wise learning rates).
    \item \textbf{IGGS-Merge}: Spatial gradient surgery with joint diagonal Fisher preconditioning.
\end{itemize}

\subsection{Main Results and Discussion}
The quantitative results are summarized in Table~\ref{tab:main_results}, and the corresponding cumulative average adaptation accuracy trajectories over test steps are illustrated in Figure~\ref{fig:cumulative_acc}. Our proposed TMS framework demonstrates a decisive victory across both Sequential and Alternating streams.

\begin{table}[t]
\caption{Test-time model merging adaptation average accuracy (\%) on Sequential and Alternating streams using MNIST, FashionMNIST, and KMNIST expert models. Bold indicates best performance.}
\label{tab:main_results}
\vskip 0.15in
\begin{center}
\begin{small}
\begin{sc}
\begin{tabular}{lcc}
\toprule
Algorithm & Sequential Acc (\%) & Alternating Acc (\%) \\
\midrule
Uniform & 27.16 & 27.16 \\
AdaMerging & 27.59 & 27.09 \\
FP-CA & 32.86 & 14.65 \\
IGGS-Merge & 21.07 & 24.83 \\
\midrule
FP-CA + TMS (Ours) & \textbf{32.94} & 14.67 \\
IGGS-Merge + TMS (Ours) & 22.73 & \textbf{26.39} \\
\bottomrule
\end{tabular}
\end{sc}
\end{small}
\end{center}
\vskip -0.1in
\end{table}

In the Sequential stream, standard AdaMerging and FP-CA suffer from significant adaptation lag at task boundaries (steps 50 and 100), where historical momentum drags the coefficients in the old direction, leading to a temporary drop in accuracy. By applying TMS, the temporal conflict is immediately detected and resolved, resulting in near-instantaneous adaptation at task boundaries and boosting the overall sequential accuracy.

In the Alternating stream, where tasks shift on every single batch, spatial gradient conflicts are severe. IGGS-Merge resolves spatial class conflicts, but suffers from temporal momentum lag across successive steps. Our unified \emph{IGGS-Merge + TMS} algorithm achieves the highest accuracy, showing that spatial and temporal gradient surgeries are complementary and crucial for robust test-time model merging in dynamic environments.

\begin{figure*}[t]
\centering
\begin{subfigure}[b]{0.48\textwidth}
    \centering
    \includegraphics[width=\textwidth]{ttmm_sequential_accuracy.png}
    \caption{Sequential Stream}
    \label{fig:seq_acc}
\end{subfigure}
\hfill
\begin{subfigure}[b]{0.48\textwidth}
    \centering
    \includegraphics[width=\textwidth]{ttmm_alternating_accuracy.png}
    \caption{Alternating Stream}
    \label{fig:alt_acc}
\end{subfigure}
\caption{Cumulative average accuracy (\%) over test steps on the Sequential and Alternating streams. TMS algorithms consistently maintain high, stable performance and adapt rapidly to task transitions (e.g., at steps 50 and 100 on the Sequential stream), whereas baseline preconditioning methods suffer from severe tracking lag or representation collapse.}
\label{fig:cumulative_acc}
\end{figure*}

\subsection{Ablation Studies and Sensitivity Analysis}
\begin{table}[t]
\caption{Ablation studies on the sensitivity floor $\tau$, test-time adaptation steps per batch $S$, and momentum decay factor $\beta$ on the Sequential and Alternating streams. Bold indicates best performance within each study block.}
\label{tab:ablation_results}
\vskip 0.15in
\begin{center}
\begin{small}
\begin{sc}
\begin{tabular}{lccccc}
\toprule
Algorithm & Floor $\tau$ & Steps $S$ & Momentum $\beta$ & Sequential Acc (\%) & Alternating Acc (\%) \\
\midrule
FP-CA & $10^{-6}$ & 1 & 0.9 & 32.86 & 14.65 \\
FP-CA + TMS & $10^{-6}$ & 1 & 0.9 & 32.94 & 14.67 \\
IGGS-Merge & $10^{-6}$ & 1 & 0.9 & 21.07 & 24.83 \\
IGGS-Merge + TMS & $10^{-6}$ & 1 & 0.9 & 22.73 & 26.39 \\
\midrule
FP-CA + TMS & 0.01 & 1 & 0.9 & 33.08 & 31.93 \\
IGGS-Merge + TMS & 0.01 & 1 & 0.9 & 34.05 & 33.69 \\
FP-CA + TMS & 0.05 & 1 & 0.9 & 32.82 & 29.40 \\
IGGS-Merge + TMS & 0.05 & 1 & 0.9 & 34.22 & 32.29 \\
FP-CA + TMS & 0.1 & 1 & 0.9 & 31.45 & 28.92 \\
IGGS-Merge + TMS & 0.1 & 1 & 0.9 & 34.40 & 31.16 \\
\midrule
FP-CA + TMS & 0.05 & 3 & 0.9 & 33.04 & 31.53 \\
IGGS-Merge + TMS & 0.05 & 3 & 0.9 & 34.02 & 33.75 \\
FP-CA + TMS & 0.05 & 5 & 0.9 & 32.70 & 32.58 \\
IGGS-Merge + TMS & 0.05 & 5 & 0.9 & 34.22 & 34.31 \\
\midrule
FP-CA + TMS & 0.01 & 1 & 0.0 & 32.21 & 29.11 \\
IGGS-Merge + TMS & 0.01 & 1 & 0.0 & 33.95 & 31.97 \\
FP-CA + TMS & 0.01 & 1 & 0.5 & 33.12 & 30.07 \\
IGGS-Merge + TMS & 0.01 & 1 & 0.5 & 33.84 & 33.03 \\
FP-CA + TMS & 0.01 & 1 & 0.9 & 33.08 & 31.93 \\
IGGS-Merge + TMS & 0.01 & 1 & 0.9 & 34.05 & 33.69 \\
FP-CA + TMS & 0.01 & 1 & 0.95 & 33.15 & 32.23 \\
IGGS-Merge + TMS & 0.01 & 1 & 0.95 & 34.06 & 33.90 \\
FP-CA + TMS & 0.01 & 1 & 0.99 & 33.03 & 32.09 \\
IGGS-Merge + TMS & 0.01 & 1 & 0.99 & 33.50 & 33.38 \\
\bottomrule
\end{tabular}
\end{sc}
\end{small}
\end{center}
\vskip -0.1in
\end{table}

To systematically analyze the behavior of our proposed Temporal Momentum Surgery (TMS) and stabilized preconditioning, we perform extensive ablation studies on the key hyperparameters: the sensitivity floor $\tau$, the number of test-time adaptation steps per batch $S$, and the momentum decay factor $\beta$. The ablation results are presented in Table~\ref{tab:ablation_results}.

\subsubsection{Effect of Sensitivity Floor $\tau$}
As discussed in Section 3.3, standard preconditioning can lead to learning rate explosion in insensitive layers. We systematically vary the sensitivity floor $\tau \in \{10^{-6}, 0.01, 0.05, 0.1\}$. 
When $\tau = 10^{-6}$ (no clamping), the learning rate of insensitive layers is unconstrained, leading to sub-optimal tracking on the Alternating stream due to quick saturation of raw coefficients. Enforcing a sensitivity floor of $\tau = 0.01$ or $\tau = 0.05$ significantly stabilizes the optimization process. This prevents gradient vanishing, allowing the model to adapt continuously and achieving a substantial boost in accuracy on the Alternating stream.

\subsubsection{Effect of Multi-Step Test-Time Adaptation}
In rapidly shifting environments, performing a single gradient update per batch may be insufficient to fully align the merging coefficients with the current task. We evaluate the effect of increasing the number of test-time adaptation steps per batch $S \in \{1, 3, 5\}$. Performing multiple steps per batch dramatically improves the tracking accuracy of our proposed methods, particularly on the challenging Alternating stream, by allowing the merging coefficients to converge closer to the optimal task-specific configuration on each step.

\subsubsection{Effect of Momentum Decay Factor $\beta$}
To investigate the role of momentum accumulation and the efficacy of our surgery scheme under varying degrees of historical dependence, we systematically vary the momentum decay factor $\beta \in \{0.0, 0.5, 0.9, 0.95, 0.99\}$ under the optimal sensitivity floor $\tau = 0.01$ and $S = 1$ step per batch.
When $\beta = 0.0$ (the no-momentum limit), our algorithm naturally simplifies to standard preconditioned gradient descent without momentum.
Increasing $\beta$ allows the optimizer to accumulate historical updates and smooth out noise in the gradient steps, which is particularly beneficial in steady environments.
However, without our proposed surgery, large values of $\beta$ would result in a massive adaptation lag during task transitions.
With TMS, our method remains robust and performs exceptionally well even at high values of $\beta$, with $\beta = 0.95$ achieving the highest overall accuracy (34.06\% on Sequential, 33.90\% on Alternating). This clearly demonstrates that TMS successfully filters out the outdated dragging components of the momentum buffer, enabling the model to reap the acceleration and stability benefits of momentum without any adaptation lag. At $\beta = 0.99$, the momentum is extremely sticky and the performance begins to slightly decay, showing that there is an optimal balance around $\beta \in [0.9, 0.95]$.

\section{Discussion and Limitations}
\label{sec:discussion}
Our proposed Temporal Momentum Surgery (TMS) is completely self-supervised and backpropagation-free for base weights, introducing negligible computational overhead (a single low-dimensional vector dot-product and scaling per step). A limitation of this work is that it assumes expert models share the same loss basin. If the experts have divergent architectures or are trained from different initializations, they cannot be merged directly in the weight space, and feature-space merging methods would be required.

Furthermore, gradient-based TTMM under self-supervised objectives like entropy minimization is vulnerable to confirmation bias and the feedback trap on highly out-of-distribution streams. Because entropy minimization encourages the model to output highly confident predictions, if the merged model initially makes a confident but incorrect prediction, gradient descent will update the merging coefficients to further reinforce this incorrect prediction. This positive feedback loop can trap the coefficients in a sub-optimal state. Combining TMS with online prototype alignment or confidence thresholding (similar to PROTO-TTMM) represents a promising direction to alleviate this issue.

\section{Conclusion}
\label{sec:conclusion}
We proposed Temporal Momentum Surgery (TMS), a mathematically principled framework to resolve momentum lag in gradient-based Test-Time Model Merging under non-stationary streams. By projecting the historical momentum onto the Riemannian normal plane of the current gradient whenever they conflict, TMS eliminates outdated gradient dragging and enables rapid, stable adaptation across both block-sequential and rapidly alternating test streams. TMS establishes a new state-of-the-art accuracy on multi-task vision streams, demonstrating the vital importance of temporal gradient management in test-time learning.

\bibliography{example_paper}
\bibliographystyle{icml2026}

%%%%%%%%%%%%%%%%%%%%%%%%%%%%%%%%%%%%%%%%%%%%%%%%%%%%%%%%%%%%%%%%%%%%%%%%%%%%%%%
%%%%%%%%%%%%%%%%%%%%%%%%%%%%%%%%%%%%%%%%%%%%%%%%%%%%%%%%%%%%%%%%%%%%%%%%%%%%%%%
% APPENDIX
%%%%%%%%%%%%%%%%%%%%%%%%%%%%%%%%%%%%%%%%%%%%%%%%%%%%%%%%%%%%%%%%%%%%%%%%%%%%%%%
%%%%%%%%%%%%%%%%%%%%%%%%%%%%%%%%%%%%%%%%%%%%%%%%%%%%%%%%%%%%%%%%%%%%%%%%%%%%%%%
\newpage
\appendix
\onecolumn
\section{Appendix: Supplementary Material}
\label{sec:appendix}

This appendix provides comprehensive details supporting the methodology, implementation, and experimental evaluations presented in the main text of our paper, \emph{Temporal Momentum Surgery for Stable and Adaptive Test-Time Model Merging}. Specifically, we include:
\begin{itemize}
    \item Full architectural and hyperparameter details for training the specialized expert models (Section~\ref{app:expert_training}).
    \item Precise specifications of the test-time adaptation streams and datasets (Section~\ref{app:stream_details}).
    \item A detailed breakdown of the test-time optimization hyperparameters (Section~\ref{app:hyperparams}).
    \item A computational complexity and execution overhead analysis of our proposed algorithms (Section~\ref{app:complexity}).
\end{itemize}

\subsection{Expert Models Training Details}
\label{app:expert_training}
To construct our multi-expert library, we fine-tune three independent expert models (MNIST-expert, FashionMNIST-expert, and KMNIST-expert) from a shared ResNet-18 base architecture pre-trained on ImageNet. The classification head of each model is adjusted to output 10 classes. The training specifications are as follows:
\begin{itemize}
    \item \textbf{Data Slices}: Each expert is trained on the first 10,000 training samples of its target dataset to simulate a data-constrained multi-task environment.
    \item \textbf{Optimizer}: We use the AdamW optimizer with a weight decay factor of $10^{-2}$.
    \item \textbf{Learning Rate}: The base learning rate is set to $10^{-4}$ with a cosine annealing schedule.
    \item \textbf{Training Epochs}: Each expert model is fine-tuned for exactly 4 epochs with a batch size of 128, which ensures high specialization on the target domain without parameter divergence from the shared loss basin.
    \item \textbf{Fisher Information Calibration}: After training, we estimate the diagonal Fisher Information elements for each model using a separate, held-out calibration set of 500 training samples. The empirical Fisher is computed by taking the average squared gradient of the log-likelihood of predictions:
    \begin{equation}
        F_p^{(k)} = \frac{1}{|D_{\text{cal}}|} \sum_{x \in D_{\text{cal}}} \left( \nabla_p \log p(y = \hat{y} \mid x; \theta_k) \right)^2
    \end{equation}
    where $\hat{y}$ is the predicted class (pseudo-label) under expert $k$, and $p$ denotes a specific parameter index in the model.
\end{itemize}

\subsection{Test-Time Stream Specifications}
\label{app:stream_details}
The evaluation datasets (MNIST, FashionMNIST, and KMNIST) are normalized using standard ImageNet mean and standard deviation parameters: $\mu = [0.485, 0.456, 0.406]$ and $\sigma = [0.229, 0.224, 0.225]$. All images are resized from their native $28 \times 28$ resolution to $224 \times 224$ pixels and replicated across 3 channels to match the ResNet-18 input structure.

We simulate non-stationary test environments by drawing batches sequentially with a fixed batch size of 64. The details of our stream designs are:
\begin{enumerate}
    \item \textbf{Sequential Stream}: Composed of 150 total batches, split into three contiguous blocks of 50 batches:
    \begin{equation}
        \text{Stream}_{\text{seq}} = \underbrace{[B_1, \dots, B_{50}]}_{\text{MNIST}} \, \Vert \, \underbrace{[B_{51}, \dots, B_{100}]}_{\text{FashionMNIST}} \, \Vert \, \underbrace{[B_{101}, \dots, B_{150}]}_{\text{KMNIST}}
    \end{equation}
    This sequence tests the ability of test-time optimization to handle sudden, block-wise distribution shifts and measure adaptation latency at task boundaries.
    \item \textbf{Alternating Stream}: Composed of 150 total batches, cycling through the three domains sequentially on a batch-by-batch basis:
    \begin{equation}
        \text{Stream}_{\text{alt}} = [B_1^{\text{MNIST}}, B_2^{\text{Fashion}}, B_3^{\text{KMNIST}}, B_4^{\text{MNIST}}, \dots, B_{150}^{\text{KMNIST}}]
    \end{equation}
    This stream provides a highly challenging setting where both spatial and temporal gradient conflicts are active on every single step, checking the stability limits of the adapters.
\end{enumerate}

\subsection{Hyperparameter Settings for Test-Time Optimization}
\label{app:hyperparams}
For all test-time model merging algorithms, we utilize consistent optimization parameters to ensure a fair comparison:
\begin{itemize}
    \item \textbf{Damping Factor ($\alpha$)}: Enforced at $\alpha = 1.0$ for all preconditioning methods.
    \item \textbf{Momentum Decay ($\beta$)}: Set to $\beta = 0.9$ for both standard SGD with momentum and our projected momentum updates.
    \item \textbf{Base Learning Rate ($\eta_0$)}: Set to $\eta_0 = 10^{-3}$ across all runs, which was selected via grid-search on a separate validation split.
    \item \textbf{Numerical Stabilizers}: The scaling parameter $\epsilon_{\text{scale}} = 10^{-6}$ and the projection stabilizer $\epsilon = 10^{-8}$.
    \item \textbf{Hyperparameter Sweeps}: In the sensitivity floor sweep (Table 2), we test $\tau \in \{10^{-6}, 0.01, 0.05, 0.1\}$. In the multi-step sweep, we evaluate $S \in \{1, 3, 5\}$ adaptation steps per batch.
\end{itemize}

\subsection{Computational and Sample Complexity Analysis}
\label{app:complexity}
A primary advantage of Test-Time Model Merging (TTMM) over standard Test-Time Adaptation (TTA) is its high computational and sample efficiency. We analyze the theoretical complexity of our proposed methods:
\begin{enumerate}
    \item \textbf{Parameter Overhead}: Unlike standard TTA methods that require updating millions of backbone weights (e.g., $11.2\times 10^6$ parameters for ResNet-18), TTMM only optimizes layer-wise merging coefficients. For ResNet-18, the number of parameter tensors is $L = 62$. With $K = 3$ experts, the total number of optimized raw parameters is $L \times K = 186$, which is negligible ($< 0.002\%$ of the backbone size). This completely prevents representation collapse and permits extremely fast gradient steps.
    \item \textbf{Computational Complexity}: The projection step in Temporal Momentum Surgery (TMS) is exceptionally lightweight. For each parameter tensor, we compute a 1D vector dot product of size $K$. The total floating-point operations (FLOPs) required for the projection across all layers is:
    \begin{equation}
        \text{FLOPs}_{\text{proj}} = \mathcal{O}(L \cdot K)
    \end{equation}
    At $L=62$ and $K=3$, this requires fewer than $600$ FLOPs per step, which is billions of times faster than a single backward pass through the ResNet backbone (which is $\approx 1.8 \times 10^9$ FLOPs). Thus, the computational overhead of TMS is completely imperceptible in practice, making it highly suitable for resource-constrained edge devices.
\end{enumerate}

\end{document}
